%% ============================================================
%%  CHAPITRE 8 : CONCLUSION ET PERSPECTIVES
%% ============================================================
\chapter{Conclusion Générale et Perspectives}
\label{chap:conclusion}

\section{Bilan du Projet}
Le projet \textbf{MASVS Audit Copilot} a été initié avec l'ambition de résoudre deux problèmes majeurs dans le domaine de la sécurité des applications mobiles : le coût temporel de l'analyse statique et le volume écrasant de faux positifs générés par les outils automatisés.

Au terme de ce travail de conception et de développement, nous avons réussi à concevoir une plateforme complète et fonctionnelle qui répond à ces défis. L'orchestration intelligente de MobSF, couplée à un moteur de traduction strict vers le nouveau standard OWASP MASVS v2, fournit une base solide pour l'évaluation de la sécurité. 

L'intégration d'un modèle de langage de grande taille (LLM) exécuté localement (Ollama/Llama 3) constitue l'innovation majeure de ce projet. Ce pipeline de triage IA a démontré sa capacité à diviser drastiquement le nombre de faux positifs, en analysant le contexte du code source de la même manière qu'un ingénieur sécurité humain, mais en une fraction du temps. De plus, l'exécution locale de l'IA garantit que le code source audité reste strictement confidentiel, répondant ainsi aux exigences de conformité les plus strictes.

Enfin, l'architecture asynchrone (FastAPI, Celery, Redis) et les interfaces développées (Dashboard Next.js moderne et CLI Typer) offrent une expérience utilisateur professionnelle, fluide et adaptée aussi bien aux auditeurs manuels qu'aux intégrations DevSecOps automatisées.

\section{Perspectives et Améliorations Futures}
Bien que la plateforme soit opérationnelle et atteigne ses objectifs principaux, plusieurs pistes d'amélioration et d'évolution peuvent être envisagées pour les versions futures :

\begin{itemize}[label=\textcolor{PrimaryBlue}{\faRocket}]
    \item \textbf{Intégration de l'Analyse Dynamique (DAST) :} Actuellement, la plateforme se concentre sur l'analyse statique (SAST). L'intégration de tests dynamiques automatisés, en exécutant l'application dans un émulateur contrôlé (par exemple via les API d'analyse dynamique de MobSF ou Frida), permettrait de détecter des vulnérabilités complexes liées à l'exécution en temps réel (manipulation mémoire, interception réseau).
    \item \textbf{Support de Modèles d'IA Spécialisés :} Tester et intégrer des modèles LLM spécifiquement entraînés sur du code source et des données de cybersécurité (comme \textit{SecureBERT} ou des modèles "Code Llama" affinés), ce qui pourrait augmenter encore la précision de la détection des faux positifs et de la génération de correctifs.
    \item \textbf{Fine-tuning du Modèle Local :} Entraîner (fine-tune) notre propre modèle Ollama sur des bases de données de vulnérabilités mobiles spécifiques et certifiées MASVS.
    \item \textbf{Extension aux Applications Multiplateformes :} Renforcer la détection et le support spécifique des frameworks hybrides modernes tels que Flutter et React Native, qui nécessitent des techniques d'analyse statique particulières (décompilation de bytecode Dart ou JavaScript).
\end{itemize}

\vspace{1cm}
\begin{center}
  \textit{La sécurité informatique est une course continue entre les concepteurs de systèmes et ceux qui cherchent à les contourner. Des outils d'assistance intelligents, tels que MASVS Audit Copilot, deviendront indispensables pour permettre aux défenseurs de garder l'avantage.}
\end{center}

\cleardoublepage
